% Bayer 彩色滤波阵列（RGGB 排列）与一个红色位置的绿色插值；宽约100mm，色值取自图像规范的低饱和度色板。
\definecolor{FormationRed}{HTML}{A7787D}
\definecolor{FormationRedFill}{HTML}{F0E5E6}
\definecolor{FormationGreen}{HTML}{829C8B}
\definecolor{FormationGreenFill}{HTML}{E8EEE9}
\definecolor{FormationBlue}{HTML}{7188A3}
\definecolor{FormationBlueFill}{HTML}{E5EAF0}
\begin{tikzpicture}[x=7mm,y=7mm,font=\figurefont\small]
  % 左：4×4 阵列，行号自上而下
  \foreach \i in {0,...,3}{
    \foreach \j in {0,...,3}{
      \pgfmathtruncatemacro{\r}{mod(\i,2)}
      \pgfmathtruncatemacro{\c}{mod(\j,2)}
      \ifnum\r=0
        \ifnum\c=0 \def\col{FormationRed}\def\fil{FormationRedFill}\def\lab{R}
        \else \def\col{FormationGreen}\def\fil{FormationGreenFill}\def\lab{G}\fi
      \else
        \ifnum\c=0 \def\col{FormationGreen}\def\fil{FormationGreenFill}\def\lab{G}
        \else \def\col{FormationBlue}\def\fil{FormationBlueFill}\def\lab{B}\fi
      \fi
      \filldraw[fill=\fil,draw=\col,line width=0.6pt] (\j,-\i) rectangle ++(1,-1);
      \node[text=BookInk] at (\j+0.5,-\i-0.5) {\lab};
    }
  }
  \draw[BookInk,line width=1pt] (0,0) rectangle (2,-2);
  \node[anchor=north,text=BookMuted,font=\figurefont\footnotesize] at (2,-4.15) {每个感光位置只记录一种颜色};
  % 右：在中心红色位置估计绿色
  \begin{scope}[xshift=45mm,yshift=-3.5mm]
    \foreach \dx/\dy/\v in {1/0/{G=60},0/-1/{G=50},2/-1/{G=70},1/-2/{G=40}}{
      \filldraw[fill=FormationGreenFill,draw=FormationGreen,line width=0.6pt] (\dx,\dy) rectangle ++(1,-1);
      \node[font=\figurefont\scriptsize] at (\dx+0.5,\dy-0.5) {\v};
    }
    \filldraw[fill=FormationRedFill,draw=FormationRed,line width=0.9pt] (1,-1) rectangle ++(1,-1);
    \node[font=\figurefont\scriptsize] at (1.5,-1.5) {R};
    \foreach \dx/\dy in {0/0,2/0,0/-2,2/-2}{
      \draw[BookRule,line width=0.6pt] (\dx,\dy) rectangle ++(1,-1);
    }
    \node[anchor=west,align=left,font=\figurefont\footnotesize] at (3.4,-1.5)
      {缺失的绿色取\\四邻平均：\\$\hat G=(60+50+70+40)/4$\\$\phantom{\hat G}=55$};
  \end{scope}
\end{tikzpicture}
